\documentclass[11pt]{article}
\usepackage{amsmath, amssymb}
\usepackage[margin=1in]{geometry}

% ============================================================================
% INTENT DEMO (2026-08-09; fallback section corrected 2026-08-26). Every
% "with intent" expression below uses the custom conversion-time commands
% (\intent, \iarg, \speakas) or a shortcut macro (\card, \determ, ...)
% understood by mathjax-bundle-intent.js.
% This file will NOT compile to PDF (the commands exist only at conversion
% time); it is for the Accessible STEM converter. Each section pairs the
% plain notation with the annotated form so before/after can be compared in
% one converted page. Expected speech: see math-conversion/intent-demo-examples.md.
%
% THE ESCAPE HATCH (and why the integrity gate does not forbid it):
% \intent{any-name}{...} accepts ANY concept name, including names in no
% vocabulary anywhere — unknown names fall back to hyphens-to-spaces plus
% "of" on the SRE channel (see phrasesFor in intent-inline.js) and are
% passed through verbatim as the MathML intent attribute, exactly as the
% MathML 4 spec intends for open-vocabulary function heads. The build's
% referential-integrity gate constrains something different: a NOTATION
% GROUP in the vocabulary may not offer a concept that has no record
% (that would make the suggest tool propose a phantom). Author-typed
% \intent names are deliberately outside the gate's jurisdiction.
% ============================================================================

\title{MathML Intent Demonstration}
\author{Accessible STEM Project}
\date{August 2026}

\begin{document}
\maketitle

\subsection*{1. One notation, many meanings: vertical bars}

Plain (a screen reader must guess --- all four read as absolute value):
the number $|x|$, the set size $|S|$, the matrix $|A|$, the vector length
$|\vec{F}|$.

Annotated (each now speaks its own meaning):
absolute value $\abs{x}$, cardinality $\card{S}$, determinant $\determ{A}$,
magnitude $\magn{\vec{F}}$.

\subsection*{2. Products that all say ``times''}

Plain: work is $W = \vec{F}\cdot\vec{d}$ and torque involves
$\vec{r}\times\vec{F}$.

Annotated: work is $W = \dotp{\vec{F}}{\vec{d}}$ and torque involves
$\crossp{\vec{r}}{\vec{F}}$.

\subsection*{3. Nonlinear dynamics}

Plain: consider $\dot{x} = 4x^2 - 16$ with fixed point $x^*$.

Annotated: consider $\intent{time-derivative}{\dot{\iarg{x}}} = 4x^2 - 16$
with fixed point $\intent{fixed-point}{{\iarg{x}}^*}$.

In optimization the same star means the optimizer:
$\intent{optimal-solution}{{\iarg{x}}^*}$.

\subsection*{4. Quantum mechanics (wrap form, multiple operands)}

Plain: the state $\left|\psi\right\rangle$, the overlap
$\left\langle\phi\middle|\psi\right\rangle$, and the matrix element
$\left\langle m\middle|A\middle|n\right\rangle$.

Annotated: the state $\intent{ket}{\left|\iarg{\psi}\right\rangle}$,
the overlap
$\intent{braket}{\left\langle\iarg{\phi}\middle|\iarg[b]{\psi}\right\rangle}$,
and the matrix element
$\intent{matrix-element}{\left\langle\iarg{m}\middle|\iarg[b]{A}\middle|\iarg[c]{n}\right\rangle}$.

\subsection*{5. Chemistry and applied notation}

Plain: at equilibrium the bracket notation $[\mathrm{HCl}]$ collides withthe interval $[0,1]$; the gradient is $\nabla f$.

Annotated: the concentration
$\intent{concentration}{\left[\iarg{\mathrm{HCl}}\right]}$ versus the closed
interval $\closedint{0}{1}$; the gradient
$\intent{gradient}{\nabla\iarg{f}}$.

\subsection*{6. Calculus shortcuts and the escape hatch}

Annotated only: the derivative $\deriv{y}{x}$, the integral
$\defint{a}{b}{f(x)}{x}$, ``n choose k'' $\binomc{n}{k}$, and a literal
override: $\speakas{the golden ratio}{\varphi}$.

\subsection*{7. Graceful fallback (concept not in any vocabulary)}

An instructor-coined name --- the order of a group, but called by a name
nobody authored: $\intent{group-size}{\left|\iarg{G}\right|}$ --- speaks
``group size of G'' via the unknown-name fallback (hyphens become spaces,
``of'' is supplied). The name used here before 2026-08-26,
\texttt{order-of-group}, no longer demonstrates fallback: the W3C open-list
bulk import made it an authored concept, so it now speaks its authored
phrase ``order of the group G'' --- a nice illustration that yesterday's
fallback can graduate into today's vocabulary without any document change.

\end{document}
